\section{The Diophantine reduction}
\label{sec:diophantine}

We combine the effective counting theorem with the size readings of
\Cref{sec:weighted}.  The distinction between their quantifiers is
essential: the counting bound is computable from the number of units,
whereas the literal numerical sizes used to exhibit a template need
only exist when the Diophantine instance has a solution.

\begin{definition}[Positive circuits]\label{n:positive-circuits}
A positive circuit is a finite arithmetic circuit with positive integer
input variables, constant $1$, and addition and multiplication gates.
A \emph{value variable} records either an input or a gate value.
A circuit equation requires equality of two designated outputs.
\end{definition}

Solvability of positive circuit equations is undecidable.  Indeed,
the Davis--Putnam--Robinson--Matiyasevich Theorem gives undecidability
of polynomial solvability over the nonnegative integers
\citep{Matiyasevich1970,Matiyasevich1971}.  Replace each nonnegative variable $x$ by
$y-1$ with $y$ positive, expand the resulting integer polynomial,
and move its positive and negative coefficient terms to opposite
sides.  Add $1$ to both sides, so each side is positive even when
one collection of terms was empty.  Positive integer coefficients
are made by sums of $1$, powers by products, and the resulting
polynomials by positive circuits.  All these transformations are
effective and preserve existence of a solution.

\begin{lemma}[Encoding a circuit by weighted comparisons]
\label{n:circuit-template}
There is an absolute stock integer $S$ with the following property.
For every positive circuit equation $\mathcal C$, one can compute
an integer $t_S(\mathcal C)$ from its syntax such that every positive
solution determines a weighted template $F$ with at most
$t_S(\mathcal C)$ units.  In every size reading of this template
supplied by \Cref{t:weighted-realization}, the actual raw sizes
give a positive solution of $\mathcal C$ in the same cell alignment.
If all cells in that reading have size at most $Q$, the resulting
solution has all its value variables at most $Q$.
\end{lemma}

\begin{proof}
For every value variable $v$, including all gate values, prescribe
a full \emph{raw cohort} $R_v$ of size $v$ per unit and a full
\emph{side cohort} $E_v$ of size $16v$ per unit.  These sizes are
literal integers when a solution is given.  The relation between
the two cohorts is recorded by a comparison of one side unit
with sixteen distinct raw units.  For a constant value $v=1$,
compare one side unit with sixteen distinct singleton markers.

For each required square side, include the primary square cohort
$Q_v$ of side $16v$, its auxiliary cohort of shape
$(4v)\times2\times2$, the full axis-reference cohorts, and all
required full offset cohorts.  A unit from $E_v$ itself may serve
as the full side reference for $Q_v$.  Every extra full reference
and every raw or side cohort receives its prescribed offset
where required; an offset cohort requires no further offset.
The formal square sides satisfy $8\mid16v$ and $16v\ge16$.
The needed full references for bounded slots are included as
well.  Distinct roles may use disjoint lots even when their
numerical sizes happen to coincide.

An addition gate $c=a+b$ is recorded by the comparison
\begin{equation}\label{n:add-comparison}
 E_c=E_a+E_b,
\end{equation}
where an equality of cohort symbols here means an equality of
the indicated totals on chosen units.  For a multiplication
gate $c=ab$, introduce a fresh value $d=a+b$, with its own raw
and side cohorts and its side-to-sixteen-raw comparison.
Require \Cref{n:add-comparison} for $d=a+b$, and use the
square comparison
\begin{equation}\label{n:product-comparison}
 Q_d=Q_a+Q_b+32E_c.
\end{equation}
Whenever needed, add the square cohorts for $a,b,d$ and their
associated roles.  The output equality is recorded using
two distinct copies of each output side cohort.  Thus even
an equality between two individual values is read as an
equality between two-unit totals.

Each comparison is between disjoint collections of a fixed
bounded number of units, and at least one collection uses
more than one unit.  In particular, \Cref{n:product-comparison}
uses one unit against two square units and thirty-two side
units, so thirty-five bounds the total support of every circuit
comparison.  Repeated occurrences of the same variable within
a comparison use distinct copies.  The pair-total reading
in \Cref{t:weighted-equal-totals} equalizes all the actual
sizes in a cohort, so these choices of copies do not change
the values being read.  The same copies may be used again
in another comparison.  The comparisons for constants
use the common marker size one.

Choose once and for all a stock $S$ large enough for
\Cref{t:weighted-product-success,t:weighted-realization},
including all bounded marker losses, all marker requests
in the readings, and the bounded number of distinct copies
needed in any one comparison.  Enlarge it by a fixed amount
if necessary for the raw/side and output comparisons.
With the comparison support and spare-copy needs fixed,
neither the number of comparisons or cohorts nor any numerical
size enters this choice.  Singleton raw units may also count
as markers.  Include the global marker-loss bound in every
cohort's stock, so enough copies survive even if all deletions
fall in one raw cohort.  The same surviving copies can be
used in different comparisons.

The role list is computable from the circuit: take its
value variables, the fresh sums for multiplication gates,
the square and auxiliary roles they require, their full
axis references, and their offsets.  Put the fixed stock
in every role and include the fixed marker stock.  There
are only finitely many prescribed roles per value or
shape, and offsets introduce no further offset recursion.
Counting these listed lots gives $t_S(\mathcal C)$.
This calculation uses no solution values.  If sizes of
different roles coincide, keeping the lots separate
still gives this same bound on their total unit count.

Now start with a positive solution and construct the
literal weighted template using its values.  All the
prescribed total equalities hold.  Apply the weighted
size readings in their one common cell alignment.  Let
$r_v$ and $e_v$ be the actual raw and side sizes; the
cohort equalization makes these well-defined positive
integers.  The raw/side and constant comparisons give
\begin{equation}\label{n:read-side}
 e_v=16r_v,\qquad r_v=1\ \hbox{for every unit constant}.
\end{equation}
In particular each side reference is even and at least
sixteen, so \Cref{t:weighted-square} applies to every
required primary square and reads its actual size as
$e_v^2$.  Addition comparisons give $r_c=r_a+r_b$.
At a multiplication gate, the fresh sum gives
$r_d=r_a+r_b$ and \Cref{n:product-comparison} becomes
\begin{equation}\label{n:factor-thirty-two}
 \bigl(16(r_a+r_b)\bigr)^2
   =(16r_a)^2+(16r_b)^2+32(16r_c).
\end{equation}
Subtracting the first two squares gives
$512r_ar_b=512r_c$, hence $r_c=r_ar_b$.
This explains the factor $32$: the cross term of
the squared side sizes is $2\cdot16^2r_ar_b$, while
one side unit for $c$ has size $16r_c$.
The doubled output comparison gives the required
output equality.  All equations concern the same
actual raw sizes, because the size readings use one
alignment simultaneously.  These sizes are themselves
cell sizes, so their bound by $Q$ proves the last
assertion.
\end{proof}

\begin{theorem}[Undecidability of finite representability]
\label{n:undecidable}
There is no algorithm which, given the tables of a finite
lattice, decides whether it is isomorphic to the congruence
lattice of a finite algebra.  Equivalently, there is no
total computable bound, in terms of lattice cardinality,
for the minimum carrier size among all representable
finite lattices.
\end{theorem}

\begin{proof}
Suppose that such a decider exists.  By
\Cref{f:cor-witness-bound}, it gives a total computable
nondecreasing function $M(N)$ that bounds a representing
carrier for every representable lattice with at most
$N$ elements.  Explicitly, enumerate the finitely many
lattice tables on at most $N$ points, retain the YES
instances using the decider, and search the finite
colored-graph witnesses for each retained table until
one is found.  The graph criterion supplies a finite
algebra on that witness carrier.  Taking the maximum
over these finitely many searches produces $M(N)$.

Fix the absolute integer $S$ from \Cref{n:circuit-template}.
Its existence is sufficient: the single integer $S$
may be hardwired in the algorithm below.  This step
does not assert an algorithm for extracting a sharp,
or even a specified, sufficient stock from the
structural arguments.  Once fixed, the function
$t_S(\mathcal C)$ is computable from the circuit syntax.
For an input circuit equation $\mathcal C$, compute
\begin{equation}\label{n:computable-witness-bound}
 \begin{aligned}
  t&=t_S(\mathcal C),&
  N&=\Phi(t),&
  m&=M(N),&
  Q&=60^m m!.
 \end{aligned}
\end{equation}

We claim that if $\mathcal C$ has any positive solution,
then it has one whose value variables are all at most
$Q$.  Given an arbitrary positive solution, use
\Cref{n:circuit-template} to exhibit its weighted
template $F$.  Its literal size parameters may be much
larger than the bounds just computed, but
\Cref{n:template-bound} still gives $|F|\le N$.
The template is representable by the coset action of
its defining ordinary subgroup interval, as in
\Cref{t:weighted-algebra}.  Hence it has a finite-algebra
representation on at most $m$ points.

Among the representations of all nontrivial lattice
images of this $F$, take one with least carrier size
$m_0$.  Since $F$ itself is one such image, $m_0\le m$.
The nontrivial-image reduction of
\Cref{t:weighted-algebra} gives an ordinary interval
for the dual of a nontrivial image, with top group
order at most
\[
 60^{m_0}m_0!\le60^m m!=Q.
\]
The product-success and size-reading
\Cref{t:weighted-realization} then provides the common
actual cell alignment, after only the allowed bounded
marker deletions.  Every read cell size is positive
and at most this top-order bound; markers have size
one.  \Cref{n:circuit-template}, with its preserved
stocks, now extracts a positive solution of
$\mathcal C$ all of whose value variables are at most
$Q$.  The side sizes are multiples of sixteen by
\Cref{n:read-side}, so every use of the square reading
in this extraction satisfies its actual numerical
hypothesis.

Thus the complete quantifier implication is
\begin{equation}\label{n:quantifier-chain}
 \begin{aligned}
 &\exists\,\mathbf v>0\quad\mathcal C(\mathbf v)
 \\
 &\quad\Longrightarrow\quad
 \exists F\quad |F|\le N,\quad F\text{ representable}
 \\
 &\quad\Longrightarrow\quad
 \exists A\quad \operatorname{Con}A\cong F,\quad |A|\le m
 \\
 &\quad\Longrightarrow\quad
 \exists[H,G]\text{ as above with }|G|\le Q
 \\
 &\quad\Longrightarrow\quad
 \exists\,\mathbf r\in\{1,\ldots,Q\}^{k}
       \quad\mathcal C(\mathbf r).
 \end{aligned}
\end{equation}
where the circuit determines $k$ and all quantities in
\Cref{n:computable-witness-bound} are computed before
any solution is known.  Auxiliary fresh sums can be
included among the $k$ value variables; restricting
back to the original values preserves the conclusion.
The reverse implication from a bounded solution to
some solution is immediate.

An algorithm may therefore compute
\Cref{n:computable-witness-bound} and search all positive
assignments up to $Q$ to the finitely many value
variables, checking each gate and the output equality.
It terminates, and \Cref{n:quantifier-chain} makes its
answer correct.  It never needs to construct a weighted
template with unknown sizes: that template was used
only to prove that an existing solution has a bounded
replacement.  This would decide positive circuit
solvability, contradicting the
Davis--Putnam--Robinson--Matiyasevich Theorem.
The equivalent carrier-bound formulation follows
again from \Cref{f:cor-witness-bound}.
\end{proof}

The structural input to this reduction is precisely the
weighted success theorem, whose exclusion arguments use
the sufficiently-general subgroup theorem in the scope
established earlier.  The reduction introduces no
additional uniformity in the numerical weights beyond
the computable counting bound proved in
\Cref{n:template-bound}.

\begin{corollary}[Undecidability of finite subgroup-interval recognition]
\label{n:interval-undecidable}
There is no deterministic Turing machine which, given the order table
of a finite nonempty lattice $L$, always halts and decides whether
$L\cong[H,G]$ for some finite group $G$ and subgroup $H\le G$, where
$[H,G]$ consists of every subgroup $K$ with $H\le K\le G$, ordered by
inclusion.  Equivalently, there is no total computable bound, in terms
of lattice cardinality, on the least possible order of the top group
$G$ among lattices admitting such a realization.
\end{corollary}

\begin{proof}
Suppose first that such a decider exists.  For each $N\ge1$, enumerate
all order-relation tables on sets of at most $N$ points, retain those
which are nonempty lattice tables, and use the decider to select the
affirmative ones.  For each selected table, enumerate finite
multiplication tables in increasing carrier size.  Test the group
axioms, enumerate the subgroups $H$, and for each $H$ enumerate every
subgroup $K$ containing it.  Testing the finitely many bijections
then decides whether the resulting full inclusion interval is
isomorphic to the selected lattice.  Each selected search terminates.
The maximum of the witness group orders found and $1$, followed by
an increasing envelope if necessary, gives a total computable
nondecreasing function $B(N)$ such that every affirmative lattice
with at most $N$ elements has a realization $[H,G]$ with $|G|\le B(N)$.
Conversely, any total computable bound on the least possible $|G|$
gives a decider by the same finite tests restricted to group orders
up to that bound.  Thus it suffices to contradict the existence of $B$.

Fix the absolute stock integer $S$ of Lemma~\ref{n:circuit-template}; as
in the preceding proof, this one integer may be hardwired.  For a
positive circuit equation $\mathcal C$, compute from its syntax
\[
 t=t_S(\mathcal C),\qquad N=\Phi(t),\qquad Q=B(N).
\]
These quantities are computed before any solution is known.
If $\mathcal C$ has a positive solution, Lemma~\ref{n:circuit-template}
exhibits a weighted template $F$ with at most $t$ units, and
\Cref{n:template-bound} gives $|F|\le N$.  This literal template is
itself a full ordinary subgroup interval by Definition~\ref{t:weighted-template}.
It is nontrivial: the circuit construction includes side units of
size $16v\ge16$.  The bound $B$ therefore supplies an ordinary
realization $F\cong[H,G]$ with $|G|\le Q$.

Use the identity $F\to F$ as the nontrivial weighted image allowed
by Convention~\ref{t:weighted-images}.  Apply the ordinary-input part of
\Cref{t:weighted-realization} to this realization with its actual
top order $M=|G|$.  It gives one common size reading whose positive
cell sizes satisfy $u_j\le M\le Q$.  By Lemma~\ref{n:circuit-template},
that reading yields a positive solution of the same equation
$\mathcal C$ with every value variable at most $Q$.

Consequently, if $\mathcal C$ has any positive solution, it has one
in $\{1,\ldots,Q\}^{k}$, where $k$ is determined by the circuit and
may include its finitely many auxiliary fresh sums.  The reverse
implication is immediate.  Computing $Q$ and testing these finitely
many assignments would decide positive-circuit solvability,
contradicting the undecidability established at the start of this
section.  The argument uses the existence of $F$ and of its reading
only to prove a bounded replacement for a solution; it does not
require either to be constructed without that solution.
\end{proof}
